\section{Discussion}

\subsection{The Performance Variation Hierarchy}

Using a single locked configuration, we observed the following hierarchy of effects on apparent Ames QSAR performance:

\begin{center}
Source heterogeneity $\gg$ Splitting method $\approx$ In-domain algorithm choice
\end{center}

The source-heterogeneity gap (scaffold CV $\to$ LDO, $\Delta = 0.390$) was an order of magnitude larger than the split-method gap ($\Delta = 0.046$), which itself was comparable to the in-domain algorithm spread (MCC range = 0.063 across all seven algorithms under fixed features and preprocessing). Three caveats apply: the hierarchy was established on a single Ames benchmark with source labels available; the LDO $\Delta = 0.390$ also reflects asymmetric test sizes and class distributions and should be read as an upper bound; and algorithm robustness was not equivalent under source shift, so algorithm choice may matter substantially more in cross-source deployment than the in-domain comparison suggests.

\subsection{What Drives the Source-Heterogeneity Gap?}

We identify four non-exclusive hypotheses, listed in roughly decreasing parsimony:

\textbf{Prior shift from prevalence imbalance.} The 19.5-fold prevalence gap (drug 58.3\% vs.\ industrial 3.0\%) alone can mechanically lower MCC: a model trained on one prevalence and evaluated at the default threshold ($t = 0.5$) on data with very different prevalence will be miscalibrated. The directional asymmetry we observe---XGBoost reaching sensitivity of only 0.062 in the industrial$\to$drug direction---is consistent with this mechanism.

\textbf{Data quality and labeling heterogeneity.} Drug-sourced compounds come from curated pharmaceutical benchmarks, while industrial compounds come from regulatory submissions (ECHA, K-REACH) where testing protocols, laboratory standards, and reporting conventions vary. Differences in assay sensitivity, strain selection, or dose ranges could shift the prevalence gap independently of any chemical difference.

\textbf{Selection bias.} Drug candidates are designed to interact with biological targets, raising the chance of incidental genotoxicity, while industrial chemicals are selected for manufacturability and stability. This is a sampling property, not a chemical-domain difference per se.

\textbf{Chemical space difference.} Pharmaceutical and industrial libraries may simply occupy different regions of property space.\cite{Lipinski_2004,Veber_2002} If this dominates, even a perfectly calibrated cross-source model would generalize poorly.

\subsection{Disentangling Prior Shift from Genuine Source Shift}

To estimate how much of the LDO collapse is attributable to prevalence imbalance alone, we evaluated each LDO direction at three thresholds: \textbf{default} ($t = 0.5$), \textbf{prior-corrected} ($t$ = training positive rate, optimal under pure prior shift), and \textbf{oracle} ($t$ maximizing MCC on the test set, an upper bound on threshold-based recalibration). If oracle MCC approaches the in-domain CV value, the LDO collapse is mostly a calibration artifact; if it remains far below, a genuine structure--activity shift exists.

For tractability, this analysis used Morgan FP-1024 plus the 13 physicochemical descriptors rather than the full 138-feature compact set used in the main pipeline. Absolute MCC values therefore differ from Table~\ref{tab:ldo}---most notably, Random Forest's industrial$\to$drug default MCC was 0.102 here versus 0.325 in the main analysis---but the internal three-threshold comparison remains valid because identical features are used at all three thresholds.

The analysis (Table~S8) yielded three findings. First, prior-shift correction provided modest improvements (mean $+0.075$ across the six configurations, range $-0.024$ to $+0.239$). Second, the oracle MCC plateaued at 0.27--0.36, well below the in-domain scaffold-CV MCC of 0.624---a residual gap of 0.26--0.35 MCC units that cannot be eliminated by any threshold tuning. Third, algorithms differed sharply in default-threshold sensitivity: Random Forest in the industrial$\to$drug direction gained $+0.243$ MCC from default to oracle (0.102 $\to$ 0.345), while gradient boosting models gained only $+0.07$ to $+0.10$.

In short, prevalence imbalance accounts for a meaningful but limited portion of the LDO collapse. A substantial residual gap remains that no recalibration can fix, and disentangling its drivers (chemical space, data quality, selection bias) would require matched-subset analysis. Regardless of cause, the practical implication is the same: published Ames QSAR performance is highly sensitive to the source composition of the evaluation set, and per-source threshold tuning alone cannot recover in-domain performance under cross-source evaluation.

\subsection{In Vivo Micronucleus and Tiered Testing}

The dissociation between ranking ability (ROC-AUC = 0.860, PR-AUC = 0.237, an 8.5-fold enrichment over the 0.028 baseline) and classification reliability (MCC CI crossing zero) has direct practical implications. For ICH M7-style binary hazard classification,\cite{ICH_M7_2017} in~vivo MN QSAR with 2D descriptors does not achieve sufficient reliability, but the enrichment over baseline suggests value as one input in tiered screening.\cite{Dobo_2012} The strong preprocessing sensitivity ($+0.22$ MCC from salt stripping) means that reported in~vivo performance depends critically on preprocessing choices that are often undocumented.\cite{Fourches_2016} Mechanistically, micronucleus induction depends on ADME, metabolic activation, tissue distribution, and DNA repair,\cite{Tweats_2007,Kirkland_2011} none of which 2D descriptors capture.

\subsection{Cross-Endpoint Overlap and Preprocessing}

The high compound overlap between endpoints ($>$90\% shared with Ames) implies that performance differences across endpoints (Ames MCC = 0.555 vs.\ in~vivo 0.311) are unlikely to come from different chemical spaces and instead reflect genuine differences in endpoint predictability from 2D structure.\cite{Tropsha_2010} Multi-task learning over all three endpoints may be a productive direction.\cite{Ramsundar_2015}

Salt stripping was endpoint-specific: beneficial for in~vivo ($+0.22$ best-model improvement), mixed for Ames ($\pm 0.03$), and generally harmful for in~vitro ($-0.03$ to $-0.13$). One plausible explanation is that some counterions are themselves clastogenically relevant in the in~vitro CA assay; a stratified analysis by salt-containing vs.\ non-salt compounds would test this directly. The QSAR curation literature\cite{Fourches_2010,Fourches_2016} emphasizes standardization, but our results suggest the optimal strategy is endpoint-dependent.

\subsection{Practical Recommendations}

We propose the following reporting standards for genotoxicity QSAR:

\begin{enumerate}
    \item Use scaffold-based splitting for primary performance reporting.
    \item When source labels are available, report cross-source LDO performance and disclose source composition with per-source positive rates.
    \item Optimize preprocessing per endpoint rather than applying uniform pipelines.
    \item Report bootstrap CIs and applicability domain coverage.
    \item For endpoints with extreme class imbalance, separately report threshold-free (ROC-AUC, PR-AUC) and threshold-dependent (MCC, balanced accuracy) metrics.
    \item When claiming algorithm superiority, test under domain shift as well as in-domain.
\end{enumerate}

\subsection{Limitations}

(a)~Source heterogeneity analysis was limited to Ames; the other endpoints lacked source annotations. (b)~Source labels reflect data provenance rather than computed chemical classification; the relative contributions of prevalence imbalance, data quality, selection bias, and chemical space difference are partially addressed by the threshold-tuning analysis (Table~S8) but cannot be fully separated without matched-subset experiments. (c)~The in~vivo learning curve used only the locked configuration. (d)~A single scaffold split was used rather than repeated holdout. (e)~The GCN did not include edge features or pre-training.\cite{Hu_2020_pretrain} (f)~No 3D descriptors, QM descriptors, or metabolite predictions were included. (g)~Within-AD vs.\ outside-AD stratification for LDO was not conducted. (h)~McNemar tests were limited to pairwise top-two comparisons; multiple-comparison correction\cite{Demsar_2006} was not applied. (i)~Table~\ref{tab:best_models} reports best observed performance across 375 configurations on the test set, which introduces model-selection bias; the locked-configuration results (Table~\ref{tab:cascade}) provide unbiased estimates.
